\documentclass[11pt]{article}

\usepackage[T1]{fontenc}
\usepackage{bookmark}
\usepackage{amsmath, amsthm, amssymb, amsfonts, mathrsfs}
\usepackage{thmtools}
\usepackage{graphicx}
\usepackage{setspace}
\usepackage{geometry}
\usepackage{float}
\usepackage{hyperref}
\usepackage{framed}
\usepackage[dvipsnames,svgnames]{xcolor}
\usepackage{tikz-cd}
\usepackage{circuitikz}
\usepackage[most]{tcolorbox}
\usepackage{enumitem}
\usepackage[center]{titlesec}
% \usepackage{mypackage} % 你的常用符号包，按需启用
\usepackage{fontawesome5}
\usetikzlibrary{positioning}

\usepackage{pgfplots}
\pgfplotsset{compat=newest}

\usepackage[backend=biber,style=alphabetic,sorting=none]{biblatex}
\addbibresource{reference.bib} % 按需启用

\tcbuselibrary{theorems}
\tcbuselibrary{breakable}

% ---------- 颜色定义 ----------
\colorlet{LightGray}{White!90!Periwinkle}
\colorlet{LightOrange}{Orange!10}
\colorlet{LightGreen}{Green!10}
\colorlet{Lightblue}{Blue!10}
\colorlet{Lightpurple}{Purple!10}
\colorlet{LightRed}{Red!10}
\colorlet{LightYellow}{Yellow!10}
\colorlet{LightCyan}{Cyan!10}
\colorlet{LightAquamarine}{Aquamarine!10}
\colorlet{LightCadetBlue}{CadetBlue!10}

\newcommand{\HRule}[1]{\rule{\linewidth}{#1}}

% ---------- 定理环境（英文） ----------
\declaretheoremstyle[
  name = Theorem,
  postheadspace={\newline},
]{thmsty}

\declaretheorem[
  style=thmsty,
  numberwithin=section,
  refname={Theorem},
]{theorem}
\tcolorboxenvironment{theorem}{
  colback=LightGray,
  breakable,
  before upper app={\setlength{\parindent}{2em}}
}

\declaretheoremstyle[name=Definition,]{thmsty}
\declaretheorem[style=thmsty,sibling=theorem,refname={Definition}]{definition}
\tcolorboxenvironment{definition}{colback=LightCyan,breakable,before upper app={\setlength{\parindent}{2em}}}

\declaretheoremstyle[name=Remark,]{thmsty}
\declaretheorem[style=thmsty,sibling=theorem,refname={Remark}]{remark}
\tcolorboxenvironment{remark}{colback=LightRed,breakable,before upper app={\setlength{\parindent}{2em}}}

\declaretheoremstyle[name=Lemma,]{thmsty}
\declaretheorem[style=thmsty,sibling=theorem,refname={Lemma}]{lemma}
\tcolorboxenvironment{lemma}{colback=Lightblue,breakable,before upper app={\setlength{\parindent}{2em}}}

\declaretheoremstyle[name=Corollary,]{thmsty}
\declaretheorem[style=thmsty,sibling=theorem,refname={Corollary}]{corollary}
\tcolorboxenvironment{corollary}{colback=Lightpurple,breakable,before upper app={\setlength{\parindent}{2em}}}

\declaretheoremstyle[name=Proposition,]{prosty}
\declaretheorem[style=prosty,sibling=theorem,refname={Proposition}]{proposition}
\tcolorboxenvironment{proposition}{colback=LightOrange,breakable,before upper app={\setlength{\parindent}{2em}}}

\declaretheoremstyle[name=Example,]{prosty}
\declaretheorem[style=prosty,sibling=theorem,refname={Example}]{example}
\tcolorboxenvironment{example}{colback=LightGreen,breakable,before upper app={\setlength{\parindent}{2em}}}

\declaretheoremstyle[name=Exercise,]{prosty}
\declaretheorem[style=prosty,numberwithin=section]{exercise}
\tcolorboxenvironment{exercise}{colback=LightAquamarine,breakable,before upper app={\setlength{\parindent}{2em}}}

\declaretheoremstyle[name=Answer,]{prosty}
\declaretheorem[style=prosty,numberwithin=section]{answer}
\tcolorboxenvironment{answer}{colback=LightCadetBlue,breakable,before upper app={\setlength{\parindent}{2em}}}

% ---------- 罗马数字 ----------
\makeatletter
\newcommand{\rmnum}[1]{\romannumeral #1}
\newcommand{\Rmnum}[1]{\expandafter\@slowromancap\romannumeral #1@}
\makeatother

% ---------- 标题格式 ----------
% article 没有 \chapter，因此删除 \titleformat{\chapter}
\titleformat{\section}{\raggedright\Large\bfseries}{\,\thesection\,}{1em}{}
\titleformat{\subsection}{\raggedright\large\bfseries}{\,\thesubsection\,}{1em}{}

\setstretch{1.2}
\geometry{
    textheight=9in,
    textwidth=5.5in,
    top=1in,
    headheight=12pt,
    headsep=25pt,
    footskip=30pt
}


% ------------------------------------------------------------------------------

% ---------- 文档信息 ----------
\title{\textbf{Seminar On Algebraic Geometry\\ Lecture Notes 1}}
\title{%
  \vspace*{-2cm}
  \includegraphics[width=\textwidth]{SKy.png}\\[1.2cm]
  \textbf{Seminar On Algebraic Geometry\\ Lecture Notes 1}%
}
\author{Zihang Zhang}
\date{Octobaer 12, 2026}

\begin{document}


\maketitle
\thispagestyle{empty}

\tableofcontents
\newpage

\section{Sheaves}

Grothendieck's creature in algebraic geometry arises mcuh from 
differential geometry. Roughly speaking, Grothendieck generalized 
both differential geometry and algebraic geometry, in order to provide 
a new geomtric view with algebraic objects. 

\begin{definition}
    Let $(X,\mathcal{T})$ be a topology space. 
    A \textbf{presheaf} $\mathscr{F}$ of  category $\mathcal{C}$ 
    is a contravarinat functor 
    $$\mathscr{F}: \mathsf{Top}(X) \to \mathscr{C},$$
    where $\mathrm{Top}(X) = (\mathcal{T}, \subseteq)$ 
    is a preorder category. 
    
    In other words, if $\mathscr{F}$ is a presheaf of $\mathrm{Ab}$, 
    then for every open subset $U$ of $X$, 
    there is an abelian group $\mathscr{F}(U)$ and for every 
    inclusion $U \subseteq V$ of open subsets of $X$, 
    there is a morphism of abelian groups 
    $\rho_{VU} = \operatorname{res}_{V,U} : \mathscr{F} (V) \to \mathscr{F} (U)$. 
    And it satisfies
    \begin{enumerate}[label=(\arabic*)]
        \item $\mathscr{F}(\varnothing) = 0$; (Not necessary)
        \item $\rho_{UU}=\operatorname{id}_{U}:\mathscr{F}(U) \to \mathscr{F}(U)$; 
        \item If $U \subseteq V \subseteq W$ are three open subsets, 
        then $\rho_{WU}=\rho_{WV}  \circ \rho_{VU}$. 
    \end{enumerate}

    \begin{center}
        \begin{tikzcd}
            U \arrow[r] \arrow[d,hook] & \mathscr{F}(U) \\
            V \arrow[r] &  \mathscr{F}(V) \arrow[u,"\rho_{VU}"']
        \end{tikzcd}
    \end{center}

    As a matter of terminology, if : $\mathscr{F}$ is a presheaf on $X$, 
    we refer to: $\mathscr{F} ( U)$ as the \textbf{sections} of 
    the presheaf $\mathscr{F}$ over the open set $U$, 
    and we sometimes use the notation $\Gamma(U,\mathscr{F})$ 
    to denote the group $\mathscr{F}(U)$. 
    We call the maps $\rho_{VU}= \operatorname{res}_{V,U}$ 
    \textbf{restriction maps}, and we sometimes 
    write $s\mid_{V}$ instead of $\rho_{VU}(s)$, 
    if $s \in \mathscr{F}(U)$. 

    We always assume that $ \mathcal{C} $ be 
    $ \mathsf{Ab} $, i.e., the category of abelian groups, 
    in order to give the category of presheaves the abelian 
    category constructure. 
\end{definition}

Simply, a presheaf is just a contravariant functor 
$ \mathscr{F} : \mathcal{D} \to \mathcal{C} $. 
A sheaf is a presheaf whose sections are, 
in a technical sense, uniquely determined by their restrictions, 
via "glueing together". 

\begin{definition}
    A presheaf $\mathscr{F}$ on a topological space $X$ 
    is a \textbf{sheaf} if it satisfies the following 
    supplementary conditions:
    \begin{enumerate}[label=(\arabic*)]
        \setcounter{enumi}{3}
        \item \textbf{Identity axiom:} 
            If $U$ is an open subset, 
            if $\{V_{i}\}$ is an open covering of $U$, 
            and if $s \in \mathscr{F}(U)$ is an element such 
            that $s \mid_{V_{i}}=0$ for all $i$, 
            then $s=0$. 
            It is equivalanet with that  
            $s \mid_{V_{i}}=t\mid_{V_{i}} \implies s=t$. 

        \item \textbf{Glueing axiom:} If $U$ is an open subset, 
        if $\{V_{i}\}$ is an open covering of $U$, 
        and if we have elements $s_{i} \in \mathscr{F}(V_{i})$ 
        for each $i$, with the property that for each $i,j$, 
        $s_{i} \mid_{V_{i}\cap V_{j}}=s_{j} \mid_{V_{i}\cap V_{j}}$ 
        then there is an (by (4), unique) element 
        $s \in \mathscr{F}(U)$ such that $s \mid_{V_{i}} =s_{i}$ 
        for each $i$. 
    \end{enumerate}
\end{definition}

\begin{example}
    \cite[Ex. II 1.1]{hartsAG}
    Let $A$ be an abelian group, and define the \textbf{constant presheaf} 
    $\underline{A}_{\mathrm{pre}}$ associated to $A$ on the 
    topological space $X$ to be the presheaf 
    $U \mapsto A$ for all $U \neq \varnothing$, 
    with restriction maps the identity. 
    In particular, the constant presheaf $\underline{\mathbb{Z}}_{\mathrm{pre}}$ 
    associated to $\mathbb{Z}$ is given by $U \mapsto \mathbb{Z}$, 
    and any section $s \in \underline{\mathbb{Z}}_{\mathrm{pre}}(U)$ 
    is an integer $s \in \mathbb{Z}$. 

    \cite[II 1.0.3]{hartsAG} 
    Let $X$ be a topological space, and $A$ an abelian group. 
    We define the \textbf{constant sheaf} $\underline{A}$ on $X$ determined 
    by $A$ as follows. 

    Give $A$ the discrete topology, and for any open set 
    $U \in X$, let $\underline{A} (U)$ be the group of all 
    continuous maps of $U$ into $A$: 
    $$\underline{A}(U):=C(U,A).$$ 
    The restriction maps are the usual restrictions 
    of maps. Since $ A $ is a discrete topological ring, 
    each $ \underline{A}(U) $ is endowed with ring structure. 
    Immediately, $ \underline{A} $ is a ring sheaf on $ X $. 

    Note that for every connected open set $U$, 
    $\underline{A} (U) \cong A$, whence the name "constant sheaf". 
    To see this, let $s \in \mathscr{A}(U)$ be a section, 
    then $U=\bigcup_{a \in A} s^{-1}(a)$. 
    Since $U$ is connected, there is only one $s^{-1}(a)=U$ 
    and the others are empty. 
    Thus, we obtain a surjective homomorphism below 
    $$\mathscr{A}(U) \to A \quad\text{with}\quad s \mapsto s(x),\, \forall x \in U. $$ 
    It is also injective, since $ s(U)=\{0\} $ if and only if 
    $ s =0 $. 

    Therefore, if $ U $ is an open set whose connected componentes 
    are open, then $ \underline{A}(U)= \prod A $, 
    one for each connected components of $ U $. 

\end{example}

Here we define sheaves of modules in order to take a 
more geometric viewpoint on sheaves. 
They are studied in \cite[II.5]{hartsAG}.

\begin{definition}
    A \textbf{ringed space} is a pair $ (X,\mathscr{O}_{X}) $ 
    consisting of a topological space $ X $ and a sheaf of rings 
    $ \mathscr{O}_{X} $ on $ X $. 
\end{definition}

\begin{definition}
    Let $ (X,\mathscr{O}_{X}) $ be a ringed space. 
    A \textbf{sheaf of $ \mathscr{O}_{X} $-modules } is a sheaf 
    $ \mathscr{F} $ on $ X $, such that for each open set $ U \subseteq X $, 
    the abelian group $ \mathscr{F}(U) $ is an $ \mathscr{O}_{X} $-module, 
    and for for each inclusion of open sets $ V \subseteq U $, 
    the restriction homomorphism $ \mathscr{F}(U) \to \mathscr{F}(V) $ 
    is compatible with the module structure via the ring homomorphism 
    $ \mathscr{O}_{X}(U) \to \mathscr{O}_X(V) $. 
\end{definition}

Far from being abstract nonsense, 
sheaf theory generalizes and unifies many geometric objects. 
Here are some examples from geometry and analysis. 

\begin{example}
Let \(X\) be a topological space. For each open subset \(U\subseteq X\), 
let \(C(U)\) denote the \(\mathbb{R}\)-algebra of continuous functions 
\(U\to \mathbb{R}\), with restriction maps given by restriction of functions. 
Then \(\mathscr{C}\), defined by \(\mathscr{C}(U)=C(U)\), is a sheaf of 
\(\mathbb{R}\)-algebras on \(X\); the glueing axiom follows from the local 
nature of continuity. In particular, if \(M\) is a smooth manifold, then the 
structure sheaf \(\mathscr{O}_{M}\) is the sheaf of smooth functions, given by
\[
\mathscr{O}_{M}(U)=C^{\infty}(U).
\]
\end{example}

\begin{example}
Let \(X\) be a complex manifold. For each open subset \(U\subseteq X\), 
let \(\mathcal{O}(U)\) denote the \(\mathbb{C}\)-algebra of holomorphic 
functions \(U\to \mathbb{C}\), with restriction maps given by restriction 
of functions. Then the structure sheaf \(\mathscr{O}_X\), defined by 
\(\mathscr{O}_X(U)=\mathcal{O}(U)\), is a sheaf of \(\mathbb{C}\)-algebras 
on \(X\): if \(\{U_i\}\) is an open cover of \(U\) and 
\(f_i\in\mathcal{O}(U_i)\) agree on all overlaps \(U_i\cap U_j\), then 
there is a unique \(f\in\mathcal{O}(U)\) restricting to each \(f_i\), 
since holomorphicity is a local property. This is the holomorphic analogue 
of the sheaf of smooth functions on a smooth manifold. 

Already here one sees why one works with sheaves rather than only global 
sections: on a compact connected complex manifold \(X\), every global 
holomorphic function is constant, so \(\mathscr{O}_X(X)\cong\mathbb{C}\), 
whereas the sheaf \(\mathscr{O}_X\) still encodes rich local holomorphic data.
        \item Let \(M\) be a smooth manifold. For each open subset \(U\subseteq M\), let
\(\Omega^k(U)\) denote the \(C^\infty(U)\)-module of smooth \(k\)-forms on
\(U\), with restriction maps given by restriction of forms. Then
\(\Omega^k\) is a sheaf of \(C^\infty_M\)-modules on \(M\). 

Crucially, \(\Omega^k\) is not merely a
sheaf of abelian groups: for \(f\in C^\infty_M(U)\) and
\(\omega\in\Omega^k(U)\), the product \(f\omega\) is defined pointwise by
\((f\omega)_p=f(p)\omega_p\), and restriction is compatible with this
structure, making \(\Omega^k\) a sheaf of modules over the structure sheaf
\(\mathscr{O}_{M}\). In particular, \(\Omega^0=\mathscr{O}_{M}\), and \(\Omega^1\) is
the sheaf of sections of the cotangent bundle \(T^*M\). 

Indeed, sheaves of modules provide an analogue of fibre bundles 
in algebraic geometry.
For example, $ \Omega^{k} $ is a $ \mathscr{O}_{M} $-module, 
and is locally free of rank $ \binom{n}{k}$ over $ \mathscr{O}_{M} $,  
hence corresponds to the vector bundle \(\bigwedge^k T^*M\), 
i.e., $ \Omega^{k} = (-,\bigwedge^k T^*M) $. 

There is a deep relation between sheaves of modules with bundles. 
For example, in differential geometry, the Swan-Serre theorem states
that: 
for $ X $ a smooth manifold with $ \mathbb{R} $-algebra of 
smooth functions $ C^{\infty}(X) $, there is an equivalence of 
categories between that of finite rank smooth vector bundles over $ X $ 
and finitely generated projective modules over $ C^{\infty}(X) $. 
See \cite[Chapter 12]{nes} for more details. 
\end{example}

\begin{example}
Let \(p:Y\to X\) be a continuous map. For each open subset \(U\subseteq X\),
define
\[
\Gamma(Y/X)(U)=\{s:U\to Y \mid p\circ s=\mathrm{id}_U,\ s\text{ continuous}\},
\]
with restriction maps given by \(s\mapsto s|_V\). Then \(\Gamma(Y/X)\) is a
sheaf on \(X\), called the \emph{sheaf of sections} of \(p\). 

If \(p: Y \to X\) is a vector bundle, then \(\Gamma(Y/X)(U)\) is a
\(C^\infty(U)\)-module, and \(\Gamma(Y/X)\) is a sheaf of
\(C^\infty_X\)-modules, locally free of rank equal to the rank of the bundle.
Indeed, the name "section" comes from such sections 
in differential geometry. 
    
\end{example}


\begin{example}
    Let $ M $ be a smooth manifold, for any open subset $ U \subseteq X $, 
    define $ C_{c}^{\infty}(U) $ be the $ \mathbb{R} $-vector space of 
    smooth functions from $ U $ to $ \mathbb{R} $ with compact support. 
    Then $ C_{c} $ forms a presheaf, which is not a sheaf in general. 

    For example, let $ M= \mathbb{R} $, 
    $ \varphi $ be a bump function such that 
    $ \operatorname{supp}(\varphi) = [-1/3,1/3]$, and 
    $ \varphi(0)=1 $. Let 
    $ f(x)= \sum_{n \in \mathbb{Z}} \varphi(x-n) $. 
    Then $ \operatorname{f}(\varphi) = \bigcup_{n \in \mathbb{Z}} [n-1/3,n+1/3] $, 
    hence $ f \notin C_{c}(\mathbb{R}) $. 
    Now let $ U_{n}:=(n-1/2,n+2/3) $ be some open subsets covering $ \mathbb{R} $. 
    In each $ U_{n} $, 
    $\operatorname{supp}(f \mid_{U_{n}})=[n-1/3,n+1/3]$ which is compact. 
    However, the global section glueing by $ f\mid_{U_{n}} $ 
    is identically $ f $, which is not in $ C_{c}^{\infty}(\mathbb{R})$. 

    An interesting result is, as the dual topological space of 
    $ C_{c}^{\infty}(U) $, the space of distributions 
    $ \mathcal{D}'(U) $ forms a sheaf. 
    For a brief introduction to this, see \cite[{4.5}]{wikiDistribution}.
\end{example}

Now we come back to classical algebraic geometry. 
In \cite[I.3]{hartsAG}, we define the ring of regular functions on 
an algebraic variety $ V $, denoted by $ \mathscr{O}(V) $. 

\begin{example}
    \cite[II.1.0.1]{hartsAG} 
    Let $ V $ be an algebraic variety over the field $ k $. 
    For each open set $ U \subseteq V $, let 
    $ \mathscr{O}_{V}(U):=\mathscr{O}(U) $ be the ring of regular 
    functions from $ U $ to $ k $. Then $ \mathscr{O}_{V} $ 
    is a presheaf of rings on $ X $. 
    
    $ \mathscr{O}_{V} $ is a sheaf, which is called 
    the \emph{sheaf of regular functions} on $ X $. 
    This is because the regularity of a function is defined locally. 
\end{example}

Now we define the morphisms between presheaves on the same space. 

\begin{definition}
    For a presheaf $\mathscr{F}$ is a contravariant functor 
    $\mathscr{F}: \mathsf{Top}(X) \to \mathsf{Ab}$, 
    it forms a functor category 
    $[\mathsf{Top}(X)^{\mathsf{op}},\mathsf{Ab}]$. 
    We define the category of presheaves as 
    $\mathsf{PSh}(X)=[\mathsf{Top}(X)^{\mathrm{op}},\mathsf{Ab}]$. 

    Thus, the morphism $\varphi: \mathscr{F} \to \mathscr{G}$ 
    is a natural transformation, 
    which consists of morphisms of abelian groups 
    $\varphi(U): \mathscr{F}(U) \to \mathscr{G}(U)$ 
    for each open set $U$, such that whenever 
    $V \subseteq U$ is an inclusion, the diagram 
    \begin{center}
        \begin{tikzcd}
\mathscr{F}(U) \arrow[r, "\varphi(U)"] \arrow[d, "\rho_{UV}"] & \mathscr{G}(U) \arrow[d, "\rho_{UV}"] \\
\mathscr{F}(V) \arrow[r, "\varphi(V)"]                        & \mathscr{G}(V)  
        \end{tikzcd}
    \end{center}
    commutes. Or we can simply say that $ \varphi(s \mid_{V})=\varphi(s)\mid_{V} $, 
    for any section $ s \in \mathscr{F}(U) $. 

    We can simply define the category of sheaves on $ X $ 
    to be the full subcategory $ \mathsf{Sh}(X) $ of $ \mathsf{PSh}(X) $. 
\end{definition}

\begin{definition}
    Let $\varphi: \mathscr{F} \to \mathscr{G}$ be a morphism of presheaves. 
    We define the presheaf kernel of $\varphi$, 
    presheaf cokernel of $\varphi$, and presheaf image 
    of $\varphi$ to be the presheaves given by 
    $U \mapsto \mathrm{ker}\,\varphi(U)$, 
    $U \mapsto \mathrm{coker}\,\varphi(U)$, 
    and $U \mapsto \mathrm{im}\,\varphi(U)$ respectively. 
\end{definition}

Indeed, $ \mathsf{PSh} $ is an abelian category. 

\begin{lemma}
    If $\varphi: \mathscr{F} \to \mathscr{G}$ is a morphism 
    of sheaves, then the presheaf kernel of $\varphi$, 
    $\mathrm{ker}\,\varphi$ is also a sheaf. 
\end{lemma}

\begin{proof}
    Notice that $ (\ker \varphi) (U)  \subseteq \mathscr{F}(U) $ 
    hence we can use the sheaf axiom on $ \mathscr{F} $. 
\end{proof}

\section{Stalks \& Germs}

Before we investigate the category of sheaves and operations on 
sheaves, we first introduct the stalks of sheaves. 

\begin{definition}[Stalks]
    If $\mathscr{F}$ is a presheaf on $X$, and if $P$ is a point 
    of $X$, we define the \emph{stalk} $\mathscr{F}_{P}$ of 
    $\mathscr{F}$ at $P$ to be the direct limit (colimit) 
    of the groups $\mathscr{F}(U)$ for all open sets $U$ 
    containing $P$, via the restriction maps 
    $\rho$, as 
    $$\mathscr{F}_{P}:= \varinjlim \mathscr{F}(U).$$
    The direct limit means that, for any open neighborhood $ U $ 
    of $ P $, there exists a homomorphism 
    $ \rho_{U} : \mathscr{F}(U) \to \mathscr{F}_{P}$, 
    such that, for compatible homomorphisms $ \lambda_{U}: \mathscr{F}(U) \to A $, 
    there exists a unique homomorphism $ \mathscr{F}_{P} \to A $. 

    Elements in $ \mathscr{F}_{P} $ are called 
    \emph{germs} of sections of $ \mathscr{F} $ at the point $ P $. 

    In fact, the index category $\mathsf{Top}_{P}(X)$ 
    with all open neighborhood $U$ of $P$ is 
    a \emph{filtered category}, since for any open neighborhood 
    $ U $ and $ V $ of $ P $, $ U \cap V $ is also an open 
    neighborhood of $ P $. 
    By \cite[1.4 D]{vakilTheRisingSea}, 
    the germ in $\mathscr{F}_{P}$ has the form of $\langle U,s \rangle$ 
    while $ U $  is the open neighborhood of $P$ and 
    $s \in \mathscr{F}(U)$ under the equivalence relation
     that 
    $$\langle U,s \rangle = \langle V,t \rangle \Longleftrightarrow 
    \,\,\exists\text{ an open neighborhood }W\subseteq U \cap V 
    \text{ of }P \text{ s.t. } s\mid_{W}=t\mid_{W}.$$ 
\end{definition}

In fact, taking stalks gives a functor $ \mathsf{Sh}(X) \to \mathsf{Ab}$, 
which is exact and left-adjoint. 

\begin{example}[Analytic germs]
    The name of germs came from analytic continuation. 
    See \cite[Chapter 8]{ahlfors}. 
\end{example}

\begin{definition}[Etale space]
    
\end{definition}

\section{Local Nature of Sheaves}

\begin{lemma}
    \label{lemma: taking stalks}
    For every $P \in U$, there is a map 
    $\mathscr{F}(U) \to \mathscr{F}_{P}$, 
    sending a section $s \in \mathscr{F}(U)$ to 
    the equivalence class of $(U,s)$, denoted $s_{P}$. 

    If $\mathscr{F}$ is a sheaf, 
    a section $s \in \mathscr{F}(U)$ of $\mathscr{F}$ over $U$ 
    is determined by its images in the stalks $\mathscr{F}_{x}$ 
    for all $x \in U$, i.e. $s=0$ if and only if 
    $s_{P}=0$, for all $P \in U$. 
    Equivalently, $s=t \in \mathscr{F}(U)$ 
    if and only if $s_{P}=t_{P} \in \mathscr{F}_{P}$, 
    for all $P \in U$.
\end{lemma}

\begin{proof}
    For any germ $ (s_{P},U_{P}) $, such that $ S_{P}=0 $, 
    we can obtain the zero section by glueing $ s_{P} $, 
    since $ U_{P} $ forms an open cover of $ X $. 
    By the identity axiom, any section glueing by $ s_{P} $ 
    must be zero. 
\end{proof}

\begin{proposition}
    \cite[2.4 N]{vakilTheRisingSea}
    Consider the morphism $\varphi:\mathscr{F} \to \mathscr{G}$ 
    of sheaves on the topology space $X$. 
    Then, the conditions below are equivalent: 
    \begin{enumerate}[label = (\arabic*)]
    \item $\varphi$ is a monomorphism. 
    \item $\varphi$ is injective on the level 
        of stalks: $\varphi_{P}:\mathscr{F}_{P}\to \mathscr{G}_{P}$ 
        is injective for all $P \in X$. 
    \item $\mathrm{ker}\,\varphi=0$.  
    \item $\varphi$ is injective on the level of open 
    sets : $\varphi(U): \mathscr{F}(U) \to \mathscr{G}(U)$ 
    is a monomorphism of $\mathrm{Ab}$, 
    thus injective and $\mathrm{ker}\,\varphi(U) = (0)$, 
    for all open set $U \subseteq X$. 
    \end{enumerate}
\end{proposition}

\begin{proof}
    We prove this by $(2) \Longrightarrow (1) \Longrightarrow (3) \Longrightarrow (4) \Longrightarrow (2)$. 

    $ (2) \implies (1) $ : By \ref{lemma: taking stalks}. 

    $ (1) \implies (3)$ : By the definition of $\mathrm{ker}\,\varphi$ 
    we could obtain an inclusion morphism 
    $i : \mathrm{ker}\,\varphi \to \mathscr{F}$ 
    and a zero morphism $0:\mathrm{ker}\,\varphi \to \mathscr{F}$ 
    sending $(\mathrm{ker}\,\varphi)(U) \mapsto (0)$ 
    inducing the commutative diagram below 
    \begin{center}
        \begin{tikzcd}
    {\mathrm{ker}\, \varphi} \arrow[r,shift left,"i"] \arrow[r,shift right,"0"'] & \mathscr{F} \arrow[r,"\varphi"] & \mathscr{G}.
    \end{tikzcd}
    \end{center}
    Since $ \varphi $ is a monomorphism, $ i $ must be $ 0 $, 
    which means that $ \ker \varphi = 0 $. 

    $ (3) \implies (4) $ : 
    Notice that $ \ker (\varphi(U))= (\ker \varphi)(U) $. 

    $ (4) \implies (2) $ :
    Consider any sheaf morphism $f,g:\mathscr{H} \to \mathscr{F}$, 
    they are given by the group homomorphisms 
    $f(U),g(U): \mathscr{H}(U) \to \mathscr{F}(U)$. 
    If $\varphi f = \varphi g$, then 
    $\varphi(U)f(U)=\varphi(U)g(U)$, 
    it follows that $f(U)=g(U)$ for $\varphi(U)$ is a 
    monomorphism in $ \mathsf{Ab} $ for any open subset 
    $U \subseteq X$. By the definition of morphisms, $f=g$. 

\end{proof}

\begin{proposition}
    \cite[2.4 O]{vakilTheRisingSea}
    Consider the morphism $\varphi:\mathscr{F} \to \mathscr{G}$ of sheaves 
    on the topology space $X$. 
    Then, the conditions below are equivalent: 
    \begin{enumerate}[label = (\arabic*)]
        \item $\varphi$ is an epimorphism. 
        \item $\varphi$ is surjective on the level of stalks: 
        $\varphi_{P}:\mathscr{F}_{P}\to \mathscr{G}_{P}$ is an 
        epimorphism, thus surjective and 
        $\mathrm{im}\,\varphi_{P} = \mathscr{G}_{P}$, 
        for all $P \in X$. 
        \item For every open set $ U \subseteq X $, 
        and for every section $ s \in \mathscr{G}(U) $, 
        there is a covering $ \{U_{i}\} $ of $ U $, 
        and sections $ t_{i} \in \mathscr{F}(U_{i}) $ 
        such that $ \varphi(t_{i})=s\mid_{U_{i}} $ for all $ i $.   
    \end{enumerate}
\end{proposition}

\begin{proof}
We prove \((1)\Rightarrow(2)\Rightarrow(3)\Rightarrow(1)\).

\((1)\Rightarrow(2)\): Assume \(\varphi\) is an epimorphism. Let
\(\operatorname{Im}\varphi\) be the image sheaf of \(\varphi\), so that
\(\varphi\) factors as
\[
\mathscr F \xrightarrow{q} \operatorname{Im}\varphi
\xrightarrow{i} \mathscr G,
\]
where \(i\) is a monomorphism. Since \(\varphi\) is an epimorphism, \(i\)
is also an epimorphism: indeed, if \(\alpha\circ i=\beta\circ i\), then
\(\alpha\circ\varphi=\beta\circ\varphi\), hence \(\alpha=\beta\). Thus
\(i\) is both a monomorphism and an epimorphism, hence an isomorphism.
Taking stalks gives
\[
\operatorname{im}\bigl(\varphi_P:\mathscr F_P\to\mathscr G_P\bigr)
=(\operatorname{Im}\varphi)_P
=\mathscr G_P,
\]
so \(\varphi_P\) is surjective for every \(P\in X\).

\((2)\Rightarrow(3)\): Let \(U\subseteq X\) be open and \(s\in\mathscr G(U)\).
For each \(P\in U\), the surjectivity of \(\varphi_P\) gives \(t_P\in
\mathscr F_P\) with \(\varphi_P(t_P)=s_P\). Choose a representative
\(t^P\in\mathscr F(V_P)\) of \(t_P\) on some open neighbourhood \(V_P\subseteq U\)
of \(P\). Then \(\varphi(t^P)\) and \(s\) have the same stalk at \(P\), so
there is an open neighbourhood \(W_P\subseteq V_P\) of \(P\) such that
\[
\varphi(t^P|_{W_P})=s|_{W_P}.
\]
The sets \(\{W_P\}_{P\in U}\) cover \(U\), and the sections
\(t^P|_{W_P}\in\mathscr F(W_P)\) satisfy the required condition.

\((3)\Rightarrow(1)\): Let \(\alpha,\beta:\mathscr G\to\mathscr H\) be morphisms
of sheaves with \(\alpha\circ\varphi=\beta\circ\varphi\). For any open
\(U\subseteq X\) and any \(s\in\mathscr G(U)\), condition (3) gives an open
cover \(\{U_i\}\) of \(U\) and sections \(t_i\in\mathscr F(U_i)\) such that
\(\varphi(t_i)=s|_{U_i}\). Then
\[
\alpha(s)|_{U_i}
=\alpha(\varphi(t_i))
=\beta(\varphi(t_i))
=\beta(s)|_{U_i}.
\]
Since the \(U_i\) cover \(U\), we get \(\alpha(s)=\beta(s)\). Hence
\(\alpha=\beta\), so \(\varphi\) is an epimorphism.

Thus \((1)\Leftrightarrow(2)\Leftrightarrow(3)\).
\end{proof}

\begin{proposition}[1.1]
    
\end{proposition}



\section{Sheafification}



\section{Historical Remark}

\printbibliography


\end{document}